\section{Incidence budgets at the affine triple-selection gate}
\label{sec:affine-common-kernel-triple-selection}

The adaptive determinant theorem leaves one precise selection problem: an
exceptional point may choose its own arm divisors, while the determinant bound
needs three points with a large \emph{common} divisor product.  We record the
finite geometry and the exact incidence budget available for this step.

Let $S\subseteq[1,M]^2\cap\mathbb Z^2$, put $N=M-1$, and retain the affine
arms $U,V,W$ of Section~\ref{sec:affine-adaptive-common-kernel}.  For a fixed
actual arm-divisor label
\[
 \lambda=(d_U,d_V,d_W),\qquad D(\lambda)=d_Ud_Vd_W,
\]
define
\[
 F_\lambda=\{x\in S:d_U\mid U(x),\ d_V\mid V(x),\ d_W\mid W(x)\}.
\]

\begin{lemma}[Integer-line cap]\label{lem:affine-integer-line-cap}
If every triple in $S$ is collinear, then $|S|\le M$.  Equivalently, more
than $M$ points in the square contain a noncollinear triple.
\end{lemma}

\begin{proof}
Choose distinct $p,q\in S$ when $|S|>1$.  If $p_1=q_1$, collinearity forces
every point of $S$ to have that first coordinate, so projection to the second
coordinate is injective and has at most $M$ values.  If $p_1\ne q_1$, and
$r,s\in S$ have $r_1=s_1$, subtracting
$\det(q-p,r-p)=\det(q-p,s-p)=0$ gives
$(q_1-p_1)(r_2-s_2)=0$.  Hence $r=s$, so projection to the first coordinate
is injective and again $|S|\le M$.
\end{proof}

\begin{theorem}[Large-label fibre and incidence bounds]
\label{thm:affine-large-label-incidence}
Assume the three arms are pairwise coprime at every point of $S$.  If
$D(\lambda)>N^2$, then $|F_\lambda|\le M$.  Consequently, for every finite
family $\mathcal L$ of such large labels,
\[
 \sum_{\lambda\in\mathcal L}|F_\lambda|\le M|\mathcal L|.
\]
If every $x$ in a set $E\subseteq S$ belongs to at least $\mu$ of these
fibres, then
\begin{equation}\label{eq:affine-large-label-multiplicity}
                         \mu|E|\le M|\mathcal L|.
\end{equation}
\end{theorem}

\begin{proof}
If $|F_\lambda|>M$, Lemma~\ref{lem:affine-integer-line-cap} produces a
noncollinear triple in the fibre.  Use the same label at all three points.
Actual arm divisibility yields the six difference congruences, and
pairwise coprimality of the arms makes $d_U,d_V,d_W$ pairwise coprime.  The
adaptive determinant theorem then gives $D(\lambda)\le N^2$, a contradiction.
Summing the fibre bound proves the first displayed incidence inequality.
Double-counting
$\{(x,\lambda):x\in E\cap F_\lambda\}$ proves
\eqref{eq:affine-large-label-multiplicity}.
\end{proof}

The associated energy identity shows exactly what an analytic refinement
must estimate.  For positive labels $d_Z(x)$, set
\[
 \mathcal E_3=\sum_{x,y,z\in S}
   \prod_{Z\in\{U,V,W\}}\gcd(d_Z(x),d_Z(y),d_Z(z))
\]
and
\[
 n(e)=\#\{x\in S:e_U\mid d_U(x),\ e_V\mid d_V(x),\ e_W\mid d_W(x)\}.
\]
Then the divisor identity $m=\sum_{e\mid m}\varphi(e)$ gives, by expanding
three finite sums,
\begin{equation}\label{eq:affine-third-gcd-energy}
 \mathcal E_3=
 \sum_{e_U,e_V,e_W}\varphi(e_U)\varphi(e_V)\varphi(e_W)n(e)^3.
\end{equation}
Thus total energy alone is insufficient: one must place enough of it on
$e_Ue_Ve_W>N^2$ and control the contribution of collinear triples.

That collinear qualification cannot be removed.  For the canonical seed
$(a,b,c)=(1,8,9)$, $R=6$ and $M=22143$.  At
\[
 (21480,282),\quad(21480,6211),\quad(21480,12140)
\]
use the common square label
\[
 (d_U,d_V,d_W)=(128881,49,121)=(359^2,7^2,11^2).
\]
The corresponding arm rows are
\[
\begin{array}{c|ccc}
 &U&V&W\\ \hline
1&128881&144109&142417\\
2&128881&464275&427009\\
3&128881&784441&711601.
\end{array}
\]
The label divides every row, all canonical admissibility and coprimality
conditions hold, and the vertical increment $5929=49\cdot121$ proves all
six congruences.  Yet
\[
 128881\cdot49\cdot121=764135449>22142^2
 \quad\text{and}\quad \det(q-p,r-p)=0.
\]
This is a full-premise counterexample to the extra implication ``common
product above the box square implies noncollinearity.''  It does not meet
the noncollinearity premise of the valid determinant theorem and therefore
does not retire the affine route.

The remaining positive gate is now quantitative.  One may close this branch
by constructing a finite family $\mathcal L$ and a multiplicity $\mu$ with
$\mu|E|>M|\mathcal L|$, or by proving that the large-product part of
\eqref{eq:affine-third-gcd-energy} exceeds every possible collinear
contribution.  No full-premise counterexample to either conditional criterion
is known.
